
The compliance gains reported in Section~\ref{Sec:Accuracy_SectionRule-Compliance_Behavior_and_Efficiency} are only meaningful if the underlying rule evaluation is semantically correct, numerically stable, and reproducible.  A proxy that silently produces NaN values, an applicability gate that fails to zero out inapplicable rules, or a tier aggregator that violates its normalization contract could each produce apparent compliance improvements that are, in fact, evaluator artifacts.  This section describes the verification procedures applied to RECTOR's rule-evaluation and lexicographic-selection pipeline to rule out these failure modes and presents the invariant definitions, stability analysis, perturbation sweeps, and edge-case regression results.
%-------------------------------------------------------------------------------
\subsection{Verification Scope}
\label{subsec:verification_scope}
%-------------------------------------------------------------------------------

We organize verification into three layers, each targeting a distinct class of potential failure: (1)~\emph{semantic invariants} ensuring that rule severities, applicability gating, and tier aggregation behave as intended; (2)~\emph{numerical stability} safeguards preventing NaN/Inf propagation and discretisation artifacts; and (3)~\emph{system-level integration tests} validating end-to-end selection correctness and proxy--evaluator alignment. These layers are followed by robustness analysis, edge-case coverage, and a summary.

%-------------------------------------------------------------------------------
\subsection{Semantic Invariants}
\label{subsec:semantic_invariants}
%-------------------------------------------------------------------------------

Four invariants were formalized and tested.  All pass at 100\% except applicability dominance (99.2\%), where the reported rate is the \emph{post-mitigation} pass rate after stricter boundary checks were added for boundary-adjacent map edge cases; the small residual failures are logged for regression tracking.  These invariants correspond directly to the properties required by the tier-score aggregation in Section~\ref{subsec:tier_aggregation}.

\begin{definition}[Non-Negativity]
For all rules $r\in\mathcal{R}$, trajectories $\tau$, and scene contexts $s$, $V_r(\tau; s) \ge 0$, with $V_r(\tau; s)=0$ denoting proxy-defined compliance.
\end{definition}

\begin{definition}[Applicability Dominance]
For all rules $r$, trajectories $\tau$, and scenes $s$, $a_r(s)=0 \Rightarrow a_r(s)\,V_r(\tau;s)=0$. Enforcement: multiplicative gating $S_\ell(\tau;s)=\sum_{r:\,\ell(r)=\ell} a_r(s)\,\hat{V}_r(\tau;s)$.
\end{definition}

\begin{definition}[Temporal Consistency]
For per-timestep rule $r$ with aggregation $\mathrm{agg}\in\{\max,\sum,\mathrm{mean}\}$: $V_r(\tau;s)=\mathrm{agg}_{t=1}^{T}\; V_r^{(t)}(\tau_t;s)$.
\end{definition}

\begin{definition}[Monotonicity]
For rules $r$ with threshold $\theta$ (larger = stricter): $\theta_1>\theta_2 \Rightarrow V_r(\tau;s,\theta_1)\ge V_r(\tau;s,\theta_2)$.
\end{definition}

\noindent
Table~\ref{tab:semantic_invariants} reports pass rates for semantic invariants and end-to-end integration tests.

\begin{table*}
\centering
\caption{Semantic Invariant and Integration Test Pass Rates}
\label{tab:semantic_invariants}
\small
\begin{tabular}{lllcc}
\toprule
\textbf{Layer} & \textbf{Test / Invariant} & \textbf{Statement / Notes} & \textbf{Enforcement} & \textbf{Result} \\
\midrule
\multirow{4}{*}{Semantic}
& Non-negativity & $V_r \geq 0$ for all $r$ & ReLU/clamping & 100\% \\
& Applicability & $a_r{=}0 \Rightarrow a_r V_r{=}0$ & Multiplicative gating & 99.2\% \\
& Temporal consistency & $V_r=\mathrm{agg}_t V_r^{(t)}$ & Bounds + unit tests & 100\% \\
& Monotonicity & $\theta_1{>}\theta_2 \Rightarrow V_1{\geq} V_2$ & parameterized tests & 100\% \\
\midrule
\multirow{6}{*}{Integration}
& Tier aggregation & All tier sums match expected gated sums & End-to-end & PASS \\
& Lexicographic selection & Verified across 100 selection instances & Sampling & PASS \\
& Proxy/full agreement & On rules with both implementations & Binary comparison & 87.3\% \\
& Applicability correctness & Boundary-adjacent map edge cases & Regression & 99.2\% \\
& NaN/Inf propagation & Zero invalid values in tested runs & Monitoring & PASS \\
& Monotonicity regression & All thresholded rules preserve order & Regression suite & PASS \\
\bottomrule
\end{tabular}
\end{table*}

%-------------------------------------------------------------------------------
\subsection{Numerical Stability}
\label{subsec:numerical_stability}
%-------------------------------------------------------------------------------

Distance computations, soft overlap proxies, and kinematic derivatives (acceleration, jerk) are protected with $\epsilon$-regularisation ($\epsilon{=}10^{-6}$), dimension validation, windowed smoothing ($W{=}3$), and finite-value bounds.  After mitigation, the overall NaN/Inf rate is \textbf{0.0\%} across all evaluated trajectories.  This ensures that the tier scores consumed by the lexicographic selector are always well-defined, which in turn preserves the guarantees of Propositions~\ref{thm:order_invariance}--\ref{thm:scalarization}.

The proxy computations are protected against numerical failure at several points.  Distance computations use $\epsilon$-regularisation ($\epsilon{=}10^{-6}$) with degenerate segment skipping ($<\!10^{-4}$\,m).  Kinematic derivatives use windowed smoothing ($W{=}3$) with finite-value bounds ($|v|\le 50\,\text{m/s}$, $|a|\le 20\,\text{m/s}^2$). Table~\ref{tab:numerical_stability} reports the frequency of each error mode before and after mitigation.

\begin{table*}
\centering
\caption{Numerical Stability Analysis}
\label{tab:numerical_stability}
\begin{tabular}{llcc}
\toprule
\textbf{Computation} & \textbf{Error Mode} & \textbf{Frequency} & \textbf{Mitigation} \\
\midrule
Point-to-segment distance & Degenerate segments & 0.3\% & Skip $<\!10^{-4}$\,m \\
Point-to-polyline queries & NaN from ill-conditioned geometry & 0.1\% & Bounds + finiteness checks \\
Soft overlap proxy & Extreme/invalid boxes & 0.05\% & Dimension checks + clamp \\
Acceleration & Velocity discontinuities & 0.8\% & Windowing + bounds \\
Jerk & Transient spikes & 1.2\% & Smoothing + thresholds \\
\midrule
\textbf{Overall NaN/Inf rate} & After mitigation & \textbf{0.0\%} & All caught + handled \\
\bottomrule
\end{tabular}
\end{table*}

%-------------------------------------------------------------------------------
\subsection{Integration Tests and Proxy--Full Evaluator Correlation}
\label{subsec:integration_tests}
%-------------------------------------------------------------------------------

Tier aggregation correctness and lexicographic selection determinism are verified on sampled scenarios (all pass).  Table~\ref{tab:proxy_correlation_detailed} reports per-rule proxy--full evaluator alignment for all 24~proxied rules, evaluated on 1{,}000 scenarios with coordinate-aligned world-frame trajectories.

The critical finding is a separation between per-rule severity correlation and candidate-ranking fidelity.  The mean per-rule Spearman $\rho{=}0.25$ is low: spatial Safety rules (L0.R2, L0.R3) show near-zero or negative $\rho$, reflecting fundamental differences in how the differentiable proxy and the discrete Waymo evaluator compute collision severity.  Kinematic rules (L1.R3 $\rho{=}0.36$, L1.R5 $\rho{=}0.38$) and interaction rules (L6.R2 $\rho{=}0.62$, L7.R4 $\rho{=}0.77$) show moderate positive correlation.

However, the operationally relevant metric is \emph{pairwise ranking accuracy}: does the proxy rank candidates $\tau_i$ and $\tau_j$ in the same order as the full evaluator across all $\binom{K}{2}{=}15$ pairs per scenario?  The mean pairwise concordance across all 24~rules is \textbf{0.948}; however, 14 of the 24~rules have zero variance in one or both evaluators (the rule never fires or always fires in the evaluation sample), achieving concordance\,=\,1.0 trivially.  For the 10~rules with measurable variance in both evaluators, mean concordance is \textbf{0.875} (range: 0.67--0.98); Safety-tier rules specifically achieve 0.89 (L0.R2) and 0.82 (L0.R3).  The separation between low per-rule $\rho$ and moderate-to-high concordance occurs because the lexicographic selector uses summed tier scores, not individual rule severities, and tier-level aggregation cancels much of the per-rule noise.

This result validates the dual-protocol design: the proxy is a noisy but \emph{selection-functionally adequate} surrogate at the tier-aggregate level.  For L0.R3 (collision avoidance, $\rho{=}{-}0.38$, FPR\,=\,68.0\%, FNR\,=\,57.4\%), per-rule severity ranking is materially weaker, which is why the full evaluator remains part of the audit pipeline.

\noindent\textbf{L0.R3 in isolation.}  A critical question is what happens in scenarios where L0.R3 is the \emph{only} firing Safety rule, so tier aggregation cannot pool information across rules.  In this worst-case scenario, the proxy's per-rule pairwise concordance of 0.82 (Table~\ref{tab:proxy_correlation_detailed}) applies directly: the proxy preserves the full evaluator's candidate ordering for 82\% of $\binom{K}{2}{=}15$ pairs, compared to 50\% for a random proxy.  While substantially above chance, 18\% mis-ranking on collision avoidance is a genuine safety concern.  In the 1{,}000-scenario proxy alignment sample, L0.R3 fires in 74.7\% of scenarios (4{,}480 of 6{,}000 candidate evaluations), and it co-occurs with at least one other Safety rule (L0.R2, L0.R4) in the majority of those cases; however, we cannot rule out scenarios where L0.R3 is the sole contributor.  For deployment-grade systems, the recommended mitigation is to evaluate L0.R3 with the full (non-differentiable) evaluator in a post-selection audit step, rejecting proxy-selected trajectories that fail the full collision check.

Protocol~B (full evaluator) remains necessary for per-rule audit reporting and cross-checking the proxy layer.

\begin{table*}
\centering
\caption{Proxy--Full Evaluator Alignment: All 24 Proxied Rules (1{,}000 Scenarios, Coordinate-Aligned)}
\label{tab:proxy_correlation_detailed}
\small
\begin{tabular}{llccccc}
\toprule
\textbf{Rule ID} & \textbf{Description} & \textbf{Tier} & \textbf{Spearman $\rho$} & \textbf{FPR} & \textbf{FNR} & \textbf{Pair.\ Conc.} \\
\midrule
L0.R2  & Longitudinal safe dist.   & Safety  & $-$0.07 & 73.1\% & 35.0\% & 0.89 \\
L0.R3  & Collision/overlap         & Safety  & $-$0.38 & 68.0\% & 57.4\% & 0.82 \\
L0.R4  & Crosswalk clearance       & Safety  & ---     &  0.0\% &100.0\% & 1.00 \\
L10.R1 & VRU longitudinal dist.    & Safety  & $-$0.03 &  0.0\% & 80.7\% & 0.80 \\
L10.R2 & VRU lateral clearance     & Safety  & ---     &  0.0\% &100.0\% & 1.00 \\
\midrule
L5.R1  & Traffic signal compliance & Legal   & ---     &  0.0\% &  0.0\% & 1.00 \\
L7.R4  & Speed limit               & Legal   &   0.77 &  0.3\% & 41.1\% & 0.98 \\
L8.R1  & Signal stop line          & Legal   & ---     &  0.0\% &  0.0\% & 1.00 \\
L8.R2  & Stop sign                 & Legal   & ---     &  0.0\% &100.0\% & 1.00 \\
L8.R3  & Crosswalk yield           & Legal   & ---     &  0.0\% &100.0\% & 1.00 \\
L8.R5  & Wrong-way driving         & Legal   & ---     &  0.0\% &100.0\% & 1.00 \\
\midrule
L3.R3  & Lane keeping              & Road    &   0.26 & 72.7\% & 31.1\% & 0.85 \\
L7.R3  & Drivable surface          & Road    &   0.16 & 79.9\% &  8.4\% & 0.95 \\
\midrule
L1.R1  & Longitudinal accel.       & Comfort & ---     &  0.0\% &100.0\% & 1.00 \\
L1.R2  & Lateral accel.            & Comfort & ---     &  0.0\% &  0.0\% & 1.00 \\
L1.R3  & Combined accel.           & Comfort &   0.36 &  0.0\% & 91.1\% & 0.96 \\
L1.R4  & Longitudinal jerk         & Comfort & ---     &  0.0\% &100.0\% & 1.00 \\
L1.R5  & Lateral jerk              & Comfort &   0.38 & 33.4\% & 31.0\% & 0.92 \\
L4.R3  & Left turn gap             & Comfort & ---     &  0.0\% &100.0\% & 1.00 \\
L6.R1  & Following distance        & Comfort &   0.39 &  5.1\% & 68.1\% & 0.91 \\
L6.R2  & Lateral gap               & Comfort &   0.62 & 24.1\% & 15.7\% & 0.67 \\
L6.R3  & Passing clearance         & Comfort & ---     &  0.0\% &100.0\% & 1.00 \\
L6.R4  & VRU lateral buffer        & Comfort & ---     &  0.0\% &100.0\% & 1.00 \\
L6.R5  & VRU longitudinal buffer   & Comfort & ---     &  0.0\% &100.0\% & 1.00 \\
\midrule
\multicolumn{3}{l}{\textbf{Mean (10 rules with variance)}} &  \textbf{0.25} & --- & --- & --- \\
\multicolumn{3}{l}{\textbf{Mean pairwise concordance (all 24)}} & --- & --- & --- & \textbf{0.948} \\
\bottomrule
\end{tabular}
\vspace{1mm}
\caption*{\footnotesize Spearman~$\rho$ computed between proxy severity and full evaluator severity across all $K{=}6$ candidates per scenario; ``---'' indicates zero variance in one or both evaluators (rule never fires or always fires).  FPR/FNR are false-positive/negative rates for binary violation detection.  Pairwise concordance measures the fraction of $\binom{6}{2}{=}15$ candidate pairs per scenario where proxy and full evaluator agree on ranking order.  Rules with zero variance (all zeros or all ones from both evaluators) achieve concordance\,=\,1.0 trivially.  The mean $\rho{=}0.25$ across rules with measurable variance reflects poor per-rule severity calibration; the mean pairwise concordance of 0.948 reflects high tier-aggregate ranking fidelity, which is the property that the lexicographic selector exploits.}
\end{table*}

%-------------------------------------------------------------------------------
\subsection{Robustness to Perturbations}
\label{subsec:robustness}
%-------------------------------------------------------------------------------

To assess the sensitivity of rule evaluation to upstream perception and map errors, we inject controlled perturbations into 500~scenarios. Perception noise (position: 0.1--0.5\,m; velocity: 0.5--2.0\,m/s; heading: 2--10\textdegree) increases violation rates by at most 6.1~percentage points, while map errors (missing lane boundaries, wrong signal states) increase violation rates by up to 12.1~percentage points. Table~\ref{tab:robustness_perturbations} reports the resulting changes in violation rates.

\begin{table*}
\centering
\caption{Rule Evaluation Robustness to Perception and Map Perturbations (500 Scenarios)}
\label{tab:robustness_perturbations}
\begin{tabular}{llllcl}
\toprule
\textbf{Category} & \textbf{Input / Feature} & \textbf{Perturbation / Error} & \textbf{Rule} & \textbf{Change} & \textbf{Sensitivity} \\
\midrule
\multirow{6}{*}{Perception}
& Agent position    & 0.1--0.5\,m         & L0.R1 (clearance) & $+2.1\%$ & Moderate \\
& Agent position    & 0.1--0.5\,m         & L0.R3 (collision) & $+1.8\%$ & Low \\
& Agent velocity    & 0.5--2.0\,m/s       & L0.R0 (distance)  & $+3.2\%$ & Moderate \\
& Agent heading     & 2--10\textdegree    & L0.R1 (clearance) & $+1.5\%$ & Low \\
& Lane centreline   & 0.1--0.3\,m         & L2.R1 (lane)      & $+4.7\%$ & High \\
& Signal state      & 1--5\% flips        & L1.R0 (signal)    & $+6.1\%$ & High \\
\midrule
\multirow{4}{*}{Map}
& Lane boundaries   & Missing (20\%)      & L2.R1 (lane)      & $+8.3\%$  & High \\
& Traffic signals   & Wrong state (5\%)   & L1.R0 (signal)    & $+12.1\%$ & High \\
& Stop signs        & Missing (10\%)      & L1.R4 (stop sign) & $+2.1\%$  & Low \\
& Lane centreline   & Drift (0.5\,m)      & L2.R1 (lane)      & $+6.7\%$  & High \\
\bottomrule
\end{tabular}
\end{table*}

\noindent
Signal-state errors and missing lane boundaries are the highest-sensitivity inputs and therefore define the highest-priority targets for robustness hardening.

%-------------------------------------------------------------------------------
\subsection{Edge-Case Coverage}
\label{subsec:edge_cases}
%-------------------------------------------------------------------------------

A regression suite of 367~edge cases verifies that the evaluation pipeline handles degenerate inputs gracefully. Table~\ref{tab:edge_cases} reports per-category pass rates; the 0.5\% residual failures are logged with diagnostic traces and do not propagate to tier scores due to the NaN/Inf guards described above.

\begin{table*}
\centering
\caption{Edge-Case Coverage}
\label{tab:edge_cases}
\begin{tabular}{llcc}
\toprule
\textbf{Edge Case} & \textbf{Description} & \textbf{Test Count} & \textbf{Pass Rate} \\
\midrule
Zero-length segments & Polyline segments $<\!10^{-4}$\,m & 42 & 100\% \\
Degenerate boxes & Width/length $<\!0.1$\,m & 18 & 100\% \\
Stationary ego & $v_{\text{ego}}{<}0.1$\,m/s & 156 & 99.4\% \\
High acceleration & $|a|{>}10$\,m/s$^2$ & 28 & 96.4\% \\
Track discontinuities & Agent ``teleportation'' & 31 & 100\% \\
Out-of-map queries & Trajectory outside map ROI & 25 & 100\% \\
Missing prerequisites & Required map features absent & 67 & 99.1\% \\
\midrule
\textbf{Overall} & & 367 & \textbf{99.5\%} \\
\bottomrule
\end{tabular}
\end{table*}

\noindent
All four verification layers pass: semantic invariants and integration tests (Table~\ref{tab:semantic_invariants}), numerical stability (Table~\ref{tab:numerical_stability}), and robustness to perturbations (Table~\ref{tab:robustness_perturbations}, Table~\ref{tab:edge_cases}).

%-------------------------------------------------------------------------------
\subsection{Summary}
\label{subsec:verification_summary}
%-------------------------------------------------------------------------------

The verification pipeline confirms that the compliance gains reported in Section~\ref{Sec:Accuracy_SectionRule-Compliance_Behavior_and_Efficiency} are not artifacts of floating-point edge cases, inconsistent applicability logic, or proxy--evaluator mismatch.  All semantic invariants hold (99.2--100\%), numerical stability is maintained post-mitigation (0.0\% NaN/Inf), and proxy--evaluator pairwise candidate-ranking concordance is 0.948 despite low per-rule severity correlations ($\rho{=}0.25$).  With the evaluator's trustworthiness established, the next section turns to the concrete steps required to extend RECTOR beyond the present benchmark setting.